\cotefile{69}
\scanpage{50}
\mspage{9}
\unsure{Convention|Convent.|commutant} : $\Delta$ objet de base

\uline{III} \uline{Relations avec} $E_6$

Soit $\Delta$ un graphe cubique

\uline{Définition} On appelle \uline{racines} de $\Delta$ les différences \inserted{$\delta'-\delta$,}
\inserted{dans $\widehat E$,} de sommets distincts non liés de $\Delta\subset\widehat E(\Delta)$

Leur ensemble $R(\Delta)$ est considéré comme partie de $E(\Delta)={}$orthogonal de $\lambda_\Delta$

$R(\Delta)$ est stable par $r\longmapsto-r$.

\deleted{\uline{Théorème} $R(\Delta)$ est un système de racines.}

\deleted{On a $(\delta'-\delta$} Si $r$ est une racine, on a

Avec $r=\delta'-\delta$.
\[
\underline{r^2=\delta'^2+\delta^2=-2}.
\]
Faisons la liste \inserted{explicite dans le cas \unsure{du|de} $\Delta_0$ \inserted{\unsure{Simplexe standart.|graphe standard}}}, --- \unsure{trouve|trouvons} que les différences \inserted{virtuelles}
$\delta'-\delta$ en question sont six à six de la façon suivante en termes de \unsure{suites|sextes|sixtes} \inserted{(``duales'')}
$(\alpha_1,\ldots\alpha_6)(\alpha'_1,\ldots,\alpha'_6)$
\deleted{$\unsure{\zeta_1\ldots\zeta_h=\zeta'_1\ldots\zeta'_h|\xi'_1-\xi_1=\xi'_2-\xi_2=}$} de sommets satisfaisant :
\[
\left\{\begin{aligned}
&\alpha_i\text{ lié à }\alpha'_j\text{ ssi }i\ne j,\quad
  \text{les }\alpha_i\text{ non liés mutuellement,}\\
&\text{les }\alpha'_i\text{ idem, en posant }\quad
 r_\alpha=\alpha'_1-\alpha_1=\alpha'_2-\alpha_2=\dots=\alpha'_6-\alpha_6
\end{aligned}\right.
\]

\deleted{\unsure{$r_0$}}

racines :
\[
\begin{aligned}
1\;&\left\{\begin{aligned}r_0&=2\eta-\textstyle\sum\zeta_i=\lambda-\eta\\
r_0&=\zeta'_{\unsure{1|a}}-\zeta_1=\zeta'_2-\zeta_2=\ldots=\zeta'_6-\zeta_6\end{aligned}\right.
&&\text{déduit de }\left\{\begin{aligned}&(\zeta_1\ldots,\zeta_6)\\&(\zeta'_1,\ldots\zeta'_{\unsure{0|6}})\end{aligned}\right.\\[4pt]
15\;&\left\{\begin{aligned}&r_{ij}=\zeta_i-\zeta_j\\&1\le i<j\le6\end{aligned}\right.
&&\text{déduits de }\left\{\begin{aligned}
&(\zeta_j,\zeta'_j,(\zeta_{i\unsure{h|k}})_{h\ne i,j},\deleted{\unsure{\zeta_{ij}}})\deleted{\unsure{\zeta'_6})}\\
&(\zeta_i,\zeta'_i,(\zeta_{jh})_{h\ne i,j})\end{aligned}\right.\\[4pt]
20\;&\left\{\begin{aligned}&r_{ijh}=\eta-\zeta_i-\zeta_j-\zeta_h\\&1\le i<j<h\le6\end{aligned}\right.
&&\text{déduits de }\left\{\begin{aligned}
&(\zeta_i,\zeta_j,\zeta_h,\zeta_{\ell m},\zeta_{mn},\zeta_{n\ell})\\
&(\zeta_{jh},\zeta_{hi},\zeta_{ij},\zeta'_n,\zeta'_\ell,\zeta'_m)\end{aligned}\right.
\end{aligned}
\]
avec \deleted{$1\le\ell<m<n\le6$} \inserted{$\ell,m,n$} les trois él\textsuperscript{ts} de $[1,6]$ distincts de $i,j,h$ (ou selon cas \unsure{prec.|p.ex.}).

et les \deleted{\unsure{différentielles|opposées}} obtenus en échangeant $\alpha$ et $\alpha'$
% archive sheet number: 49
